% frontmatter/nivel2/detective.tex -- para la niña, no para el adulto:
% cómo se lee "con lupa" (las normas del club, que aparecen en la
% semana 3) y el carné del Club de la Lupa para rellenar el primer día.
\section*{Cómo lee un detective}

{\large
\begin{enumerate}[leftmargin=9mm, itemsep=3mm, label=\textbf{\arabic*.}]
\item \textbf{Lee el texto entero}, sin prisa, para saber de qué va.
\item \textbf{Léelo otra vez}, y ahora busca los detalles: quién, dónde,
  cuándo, cuántos, de qué color, por qué.
\item Si una frase es muy larga, \textbf{busca sus partes}: lo que pasa,
  y las palabras que lo explican -- \textit{que}, \textit{porque},
  \textit{cuando}, \textit{aunque}, \textit{si}, \textit{para que}.
\item \textbf{Antes de responder, busca la prueba}: la frase del texto que
  te da la razón.
\item Si el caso dura toda la semana, \textbf{vuelve a las páginas de los
  días anteriores}: las pistas no se escapan, se quedan ahí.
\end{enumerate}
}

\vspace{8mm}
\begin{center}
\begin{tcolorbox}[enhanced, width=0.86\linewidth, colback=white,
  colframe=colorLectura, boxrule=1.4pt, arc=3mm,
  left=5mm, right=5mm, top=4mm, bottom=5mm,
  title={\large Club de la Lupa \textperiodcentered\ Carné de detective},
  colbacktitle=colorLectura, coltitle=white, fonttitle=\bfseries\sffamily]
\begin{tabularx}{\linewidth}{@{}X r@{}}
\begin{minipage}[t]{0.56\linewidth}
\vspace{4mm}
\large
Nombre: \hrulefill\par\vspace{7mm}
Edad: \hrulefill\par\vspace{7mm}
Me fijo mucho en: \hrulefill\par\vspace{7mm}
Firma: \hrulefill
\end{minipage}
&
\begin{minipage}[t]{0.36\linewidth}
\vspace{2mm}
\centering
\framebox[40mm]{\rule{0pt}{44mm}}\par\vspace{1mm}
{\footnotesize\color{colorGris} Tu retrato, dibujado por ti}
\end{minipage}
\end{tabularx}
\end{tcolorbox}

\vspace{6mm}
\iconoLupa[0.45]
\end{center}
\newpage
